% ============================================================
%  DISCOVERY CURRICULUM — Journey 3, Module 2
%  "Matrices and Systems"
%  Compile with: pdflatex J3_M2.tex  (run twice for references)
% ============================================================

\documentclass[12pt, a4paper]{article}

\usepackage[a4paper, top=2.5cm, bottom=2.5cm, left=2.8cm, right=2.8cm, headheight=15pt]{geometry}
\usepackage{amsmath, amssymb, amsthm}
\usepackage{xcolor}
\usepackage{mdframed}
\usepackage{enumitem}
\usepackage{tikz}
\usetikzlibrary{arrows.meta, positioning, calc}
\usepackage{booktabs}
\usepackage{fancyhdr}
\usepackage{titlesec}
\usepackage{parskip}
\usepackage{hyperref}
\hypersetup{colorlinks=true, linkcolor=NavyBlue, urlcolor=NavyBlue}

% ── Colours ─────────────────────────────────────────────────
\definecolor{NavyBlue}{HTML}{1A3A5C}
\definecolor{MidBlue}{HTML}{2E6DA4}
\definecolor{Gold}{HTML}{C8922A}
\definecolor{LightBlue}{HTML}{EBF2FA}
\definecolor{LightGold}{HTML}{FDF4E3}
\definecolor{DarkGrey}{HTML}{3A3A3A}
\definecolor{BoxGreen}{HTML}{E8F5E9}
\definecolor{SidebarGrey}{HTML}{F5F5F5}

\titleformat{\section}{\Large\bfseries\color{NavyBlue}}{}{0em}{}[{\color{MidBlue}\vspace{2pt}\rule{\linewidth}{1.5pt}}]
\titleformat{\subsection}{\large\bfseries\color{MidBlue}}{}{0em}{}

\pagestyle{fancy}
\fancyhf{}
\fancyhead[L]{\small\color{DarkGrey}\textit{Discovery Curriculum}}
\fancyhead[R]{\small\color{DarkGrey}\textit{Journey 3 $\cdot$ Module 2}}
\fancyfoot[C]{\small\color{MidBlue}\thepage}
\renewcommand{\headrulewidth}{0.4pt}

\newmdenv[backgroundcolor=NavyBlue!7,linecolor=NavyBlue,linewidth=2pt,leftline=true,rightline=false,topline=false,bottomline=false,innerleftmargin=14pt,innerrightmargin=10pt,innertopmargin=10pt,innerbottommargin=10pt,skipabove=14pt,skipbelow=10pt]{puzzlebox}
\newmdenv[backgroundcolor=LightGold,linecolor=Gold,linewidth=2pt,leftline=true,rightline=false,topline=false,bottomline=false,innerleftmargin=14pt,innerrightmargin=10pt,innertopmargin=10pt,innerbottommargin=10pt,skipabove=14pt,skipbelow=10pt]{discoverybox}
\newmdenv[backgroundcolor=LightBlue,linecolor=MidBlue,linewidth=1.5pt,roundcorner=4pt,innerleftmargin=14pt,innerrightmargin=10pt,innertopmargin=10pt,innerbottommargin=10pt,skipabove=14pt,skipbelow=10pt]{namebox}
\newmdenv[backgroundcolor=BoxGreen,linecolor=DarkGrey!35,linewidth=1pt,roundcorner=4pt,innerleftmargin=14pt,innerrightmargin=10pt,innertopmargin=10pt,innerbottommargin=10pt,skipabove=14pt,skipbelow=10pt]{reflectionbox}
\newmdenv[backgroundcolor=SidebarGrey,linecolor=Gold,linewidth=1.2pt,leftline=true,rightline=false,topline=false,bottomline=false,innerleftmargin=12pt,innerrightmargin=10pt,innertopmargin=7pt,innerbottommargin=7pt,skipabove=10pt,skipbelow=10pt]{thought}

\newcounter{problem}
\newenvironment{problem}[1][]{\refstepcounter{problem}\par\vspace{6pt}\noindent\textbf{\color{MidBlue}Problem \theproblem}\ifx&#1&\else\ \textit{(#1)}\fi\textbf{.}\enskip}{\par\vspace{4pt}}
\newcounter{deepproblem}
\newenvironment{deepproblem}[1][]{\refstepcounter{deepproblem}\par\vspace{6pt}\noindent\textbf{\color{NavyBlue}$\bigstar$\ Deep Problem \thedeepproblem}\ifx&#1&\else\ \textit{(#1)}\fi\textbf{.}\enskip}{\par\vspace{4pt}}

\begin{document}

\thispagestyle{empty}
\begin{center}
{\small\color{MidBlue}\textbf{DISCOVERY CURRICULUM}\quad$\cdot$\quad\textbf{JOURNEY 3: LINEAR STRUCTURE}}
\vspace{0.6cm}

{\color{NavyBlue}\rule{\linewidth}{2pt}}
\vspace{0.5cm}

{\Huge\bfseries\color{NavyBlue}Module 2}\\[0.3cm]
{\LARGE\bfseries\color{MidBlue}Matrices and Systems}\\[0.4cm]
{\color{Gold}\rule{0.4\linewidth}{1.2pt}}\\[0.4cm]
{\large\itshape\color{DarkGrey}
A matrix is a linear map wearing coordinates.\\[0.15cm]
Row reduction is the act of finding its soul.}

\vspace{1cm}
\begin{tabular}{rl}
\color{DarkGrey}\textbf{Journey}     & Journey 3 --- Linear Structure \\[4pt]
\color{DarkGrey}\textbf{Module}      & 2 of 6 \\[4pt]
\color{DarkGrey}\textbf{Prerequisites} & J3 Module 1; basic matrix notation \\[4pt]
\color{DarkGrey}\textbf{You will discover} & Matrix multiplication, row reduction, determinants, invertibility \\[4pt]
\color{DarkGrey}\textbf{Format}      & 8 problems $+$ 3 deep problems $\approx$ 13 pages \\
\end{tabular}
\vspace{1cm}

{\color{NavyBlue}\rule{\linewidth}{2pt}}
\end{center}

\vspace{0.5cm}
\noindent\textbf{From abstraction to computation.}
Module~1 showed that a linear map is completely determined by its action on a basis. Once bases are chosen, a linear map $T : \mathbb{R}^n \to \mathbb{R}^m$ becomes a matrix $A$ --- an $m \times n$ array of numbers. All the abstract theory (kernel, image, rank-nullity) translates into questions about matrices that can be answered by computation. This module builds that computational engine.

\newpage

\section{Part 1 \texorpdfstring{---}{--} The Puzzle}

\begin{puzzlebox}
\textbf{The Opening Puzzle.}

Three factories produce three goods. Each factory uses the goods as inputs in the following amounts per unit produced:

\begin{center}
\begin{tabular}{lccc}
\toprule
               & Good 1 & Good 2 & Good 3 \\
\midrule
Factory 1 uses & 0.2    & 0.0    & 0.1    \\
Factory 2 uses & 0.4    & 0.3    & 0.1    \\
Factory 3 uses & 0.1    & 0.2    & 0.3    \\
\bottomrule
\end{tabular}
\end{center}

If the economy demands outputs $\mathbf{d} = (100, 200, 150)$ units, and the total production vector is $\mathbf{x}$, then the production satisfies:
\[
  \mathbf{x} - A\mathbf{x} = \mathbf{d} \qquad \Longrightarrow \qquad (I - A)\mathbf{x} = \mathbf{d}.
\]

\textbf{Question 1.} What does the system $(I-A)\mathbf{x} = \mathbf{d}$ mean? What are we solving for?

\textbf{Question 2.} This is a \emph{linear system}: three equations in three unknowns. How would you solve it?

\textbf{Question 3.} What would it mean for this system to have no solution? More than one solution? Exactly one? Which outcome do you expect here, and why?
\end{puzzlebox}

\medskip
This is the \textbf{Leontief input-output model}, developed by Wassily Leontief in the 1940s to model national economies. It won him the Nobel Prize in Economics in 1973. The mathematics is a linear system, and its solvability depends entirely on properties of the matrix $(I-A)$.

\section{Part 2 \texorpdfstring{---}{--} Exploration}

\subsection{Step 1: Matrix multiplication as composition}

\begin{problem}
\label{prob:matmul}
\begin{enumerate}[label=(\alph*), topsep=2pt]
  \item If $A$ is the matrix of $T : \mathbb{R}^n \to \mathbb{R}^m$ and $B$ is the matrix of $S : \mathbb{R}^p \to \mathbb{R}^n$, the composition $T \circ S$ has matrix $AB$. Write out the formula for $(AB)_{ij}$.
  \item Compute $AB$ and $BA$ for $A = \begin{pmatrix}1&2\\3&4\end{pmatrix}$, $B = \begin{pmatrix}0&1\\1&0\end{pmatrix}$. Is $AB = BA$?
  \item Prove that matrix multiplication is associative: $(AB)C = A(BC)$.
  \item The identity matrix $I$ satisfies $AI = IA = A$. Write $I$ explicitly for $n=3$ and verify.
  \item Geometric interpretation: if $A$ represents a $45°$ rotation and $B$ represents a scaling by 2, what does $AB$ represent? Does the order matter?
\end{enumerate}
\end{problem}

\begin{thought}
\textbf{Why $AB \neq BA$ in general.}
Matrix multiplication represents composition of linear maps. And composition of functions is generally not commutative: rotating then reflecting is not the same as reflecting then rotating. Non-commutativity is not a flaw in matrices --- it is an accurate reflection of the geometry. This non-commutativity is also why Heisenberg's uncertainty principle holds: the operators for position and momentum do not commute, and their commutator is proportional to Planck's constant. The algebra of matrices is, in a deep sense, the algebra of quantum mechanics.
\end{thought}

\subsection{Step 2: Row reduction and linear systems}

\begin{problem}
\label{prob:rowred}
The system $A\mathbf{x} = \mathbf{b}$ is solved by row-reducing the augmented matrix $[A \mid \mathbf{b}]$ to row echelon form.
\begin{enumerate}[label=(\alph*), topsep=2pt]
  \item Row reduce and solve:
  $\begin{pmatrix}1&2&3\\2&5&7\\1&1&2\end{pmatrix}\mathbf{x} = \begin{pmatrix}6\\13\\5\end{pmatrix}$

  \item Row reduce to find all solutions:
  $\begin{pmatrix}1&2&1&3\\2&4&3&7\\1&2&0&2\end{pmatrix}\mathbf{x} = \begin{pmatrix}0\\0\\0\end{pmatrix}$

  \item The system $A\mathbf{x} = \mathbf{0}$ always has the trivial solution $\mathbf{x} = \mathbf{0}$. When does it have other solutions? Express your answer in terms of rank and nullity.

  \item A system $A\mathbf{x} = \mathbf{b}$ has: no solution if $\mathbf{b} \notin \text{im}\,A$; a unique solution if $\ker A = \{\mathbf{0}\}$; infinitely many if $\ker A \neq \{\mathbf{0}\}$. Verify all three cases with examples.

  \item Solve the Leontief system from the puzzle: $(I-A)\mathbf{x} = \mathbf{d}$ with the given matrix and $\mathbf{d} = (100, 200, 150)^T$.
\end{enumerate}
\end{problem}

\begin{problem}
\label{prob:rref}
The \textbf{reduced row echelon form} (RREF) of a matrix encodes all structural information.
\begin{enumerate}[label=(\alph*), topsep=2pt]
  \item Find the RREF of $A = \begin{pmatrix}2&4&1&3\\0&0&3&6\\1&2&2&6\end{pmatrix}$.
  \item Identify the pivot columns and free columns. The pivot columns form a basis for the column space (image). What is $\text{rank}(A)$?
  \item Write the general solution to $A\mathbf{x} = \mathbf{0}$ in terms of free variables.
  \item Verify rank-nullity: $\text{rank}(A) + \text{nullity}(A) = $ number of columns.
  \item The RREF is unique (same for any sequence of row operations). Why does this make it a canonical form for the row space?
\end{enumerate}
\end{problem}

\subsection{Step 3: Invertibility and determinants}

\begin{problem}
\label{prob:invert}
A square matrix $A$ is \textbf{invertible} if there exists $A^{-1}$ with $AA^{-1} = A^{-1}A = I$.
\begin{enumerate}[label=(\alph*), topsep=2pt]
  \item $A$ is invertible if and only if $A\mathbf{x} = \mathbf{0}$ has only the trivial solution. Prove this.
  \item Find $A^{-1}$ for $A = \begin{pmatrix}2&1\\5&3\end{pmatrix}$ by row-reducing $[A \mid I]$ to $[I \mid A^{-1}]$.
  \item Prove: $(AB)^{-1} = B^{-1}A^{-1}$. The order reversal matters --- explain geometrically.
  \item The \textbf{determinant} $\det A$ measures how much $A$ scales volumes. For $2\times 2$: $\det\begin{pmatrix}a&b\\c&d\end{pmatrix} = ad - bc$. Compute $\det A$ for the matrix in (b). What is the connection between $\det A = 0$ and non-invertibility?
  \item For a $3\times 3$ matrix, expand the determinant along the first row using cofactors. Compute $\det\begin{pmatrix}1&2&3\\0&1&4\\5&6&0\end{pmatrix}$.
\end{enumerate}
\end{problem}

\begin{problem}
\label{prob:det_props}
Determinants have three defining properties from which everything else follows.
\begin{enumerate}[label=(\alph*), topsep=2pt]
  \item The three properties: (i) $\det I = 1$; (ii) $\det$ changes sign when two rows are swapped; (iii) $\det$ is linear in each row separately. Use these to derive:
    \begin{itemize}[noitemsep]
      \item $\det A = 0$ if two rows are equal
      \item $\det(cA) = c^n \det A$ for an $n\times n$ matrix
      \item $\det(AB) = \det A \cdot \det B$
    \end{itemize}
  \item Verify $\det(AB) = \det A \cdot \det B$ numerically for two $2\times 2$ matrices.
  \item A matrix representing a rotation by angle $\theta$ in $\mathbb{R}^2$ is $R = \begin{pmatrix}\cos\theta & -\sin\theta \\ \sin\theta & \cos\theta\end{pmatrix}$. Compute $\det R$. Does a rotation change volumes?
  \item A \textbf{reflection} matrix has $\det = -1$. What does the sign of the determinant represent geometrically?
  \item \textbf{Cramer's rule}: the solution to $A\mathbf{x} = \mathbf{b}$ (when $\det A \neq 0$) has $x_i = \det A_i / \det A$, where $A_i$ is $A$ with the $i$-th column replaced by $\mathbf{b}$. Verify for the $2\times 2$ case.
\end{enumerate}
\end{problem}

\subsection{Step 4: Column space, row space, null space}

\begin{problem}
\label{prob:four_subspaces}
Every matrix $A$ ($m\times n$) has four fundamental subspaces:
\[
  \text{Column space } C(A) \subseteq \mathbb{R}^m, \quad
  \text{Row space } C(A^T) \subseteq \mathbb{R}^n, \quad
  \text{Null space } N(A) \subseteq \mathbb{R}^n, \quad
  \text{Left null space } N(A^T) \subseteq \mathbb{R}^m.
\]
\begin{enumerate}[label=(\alph*), topsep=2pt]
  \item For $A = \begin{pmatrix}1&2&3\\4&5&6\\7&8&9\end{pmatrix}$: find a basis for each of the four fundamental subspaces.
  \item What are the dimensions of each subspace? Verify they satisfy the rank-nullity theorem in both directions.
  \item \textbf{Fundamental theorem of linear algebra}: $N(A) \perp C(A^T)$ and $N(A^T) \perp C(A)$. Verify orthogonality: show that any vector in $N(A)$ is orthogonal to any vector in $C(A^T)$.
  \item $A\mathbf{x} = \mathbf{b}$ has a solution if and only if $\mathbf{b} \perp N(A^T)$ (i.e.\ $\mathbf{b}$ is orthogonal to the left null space). Verify for a specific example.
\end{enumerate}
\end{problem}

\begin{problem}
\label{prob:change_of_basis}
Suppose the same linear map $T$ is represented by matrix $A$ in basis $\mathcal{B}$ and by matrix $B$ in basis $\mathcal{C}$. Let $P$ be the change-of-basis matrix from $\mathcal{C}$ to $\mathcal{B}$.
\begin{enumerate}[label=(\alph*), topsep=2pt]
  \item Show $B = P^{-1}AP$. Matrices related this way are called \textbf{similar}.
  \item Compute $P^{-1}AP$ for $A = \begin{pmatrix}3&1\\0&2\end{pmatrix}$ and $P = \begin{pmatrix}1&1\\1&-1\end{pmatrix}$.
  \item Two similar matrices have the same determinant, the same trace ($\text{tr}\,A = \sum a_{ii}$), and the same characteristic polynomial. Verify these for $A$ and $B = P^{-1}AP$ from (b).
  \item The trace and determinant are \emph{invariants} of the linear map --- independent of the basis. List two more invariants. Why do invariants matter?
  \item \textit{(Preview.)} Can we always choose a basis in which the matrix is diagonal? When is this possible? This is the eigenvalue question, which is Module 3.
\end{enumerate}
\end{problem}

\section{Part 3 \texorpdfstring{---}{--} Discovery}

\begin{discoverybox}
\textbf{What have we built?}

\textbf{First.} Matrix multiplication is composition of linear maps. It is associative but generally not commutative --- and both properties are geometrically meaningful.

\textbf{Second.} Row reduction is the computational workhorse: it solves linear systems, computes rank, finds bases for all four fundamental subspaces, and determines invertibility.

\textbf{Third.} The determinant measures signed volume scaling. It vanishes exactly when the matrix is singular (non-invertible), and it satisfies $\det(AB) = \det A \cdot \det B$.

\textbf{Fourth.} The four fundamental subspaces --- column space, row space, null space, left null space --- and their orthogonality relations are the geometric heart of the theory. Every solvability question about $A\mathbf{x} = \mathbf{b}$ is answered by these subspaces.

\medskip
\textbf{The key question not yet answered.} Change of basis produces similar matrices. Is there a ``best'' basis --- one in which the matrix is diagonal or nearly so? The answer involves eigenvalues. When a matrix is diagonal, it acts on each basis vector by simply scaling it. Finding those basis vectors is the eigenvalue problem, the deepest and most powerful idea in linear algebra.
\end{discoverybox}

\section{Part 4 \texorpdfstring{---}{--} Naming the Ideas}

\begin{namebox}
\textbf{Matrix operations and rank.}

$(AB)_{ij} = \sum_k A_{ik}B_{kj}$ (row of $A$ dots column of $B$).

$\text{rank}(A) = \dim C(A) = \dim C(A^T)$. \quad $\text{rank}(A) + \text{nullity}(A) = n$ (columns).

$A$ ($n\times n$) is invertible $\iff$ $\det A \neq 0$ $\iff$ $\text{rank}(A) = n$ $\iff$ $A\mathbf{x}=\mathbf{0}$ has only $\mathbf{x}=\mathbf{0}$.
\end{namebox}

\begin{namebox}
\textbf{Determinant.}

Uniquely defined by: $\det I = 1$; sign-change under row swap; linearity in each row.

Key facts: $\det(AB) = \det A \cdot \det B$; $\det(A^T) = \det A$; $\det(A^{-1}) = 1/\det A$.

$|\det A|$ = volume scaling factor of the linear map represented by $A$.
\end{namebox}

\section{Part 5 \texorpdfstring{---}{--} Deep Problems}

\begin{deepproblem}[LU decomposition]
Every invertible matrix $A$ can be written as $A = LU$ where $L$ is lower triangular with ones on the diagonal and $U$ is upper triangular.
\begin{enumerate}[label=(\alph*)]
  \item Perform Gaussian elimination on $A = \begin{pmatrix}2&1&1\\4&3&3\\8&7&9\end{pmatrix}$ keeping track of the multipliers used. Show that the multipliers form the matrix $L$.
  \item Why is $LU$ decomposition faster than recomputing Gaussian elimination for each right-hand side $\mathbf{b}$ in $A\mathbf{x}=\mathbf{b}$?
  \item When row swaps are needed, the decomposition becomes $PA = LU$ where $P$ is a permutation matrix. Show $\det P = \pm 1$ and explain the sign.
  \item The number of floating-point operations for Gaussian elimination on an $n\times n$ matrix is $O(n^3)$. For $n = 1000$, how does this compare to $O(n^2)$ algorithms? Why does this matter for large-scale scientific computing?
\end{enumerate}
\end{deepproblem}

\begin{deepproblem}[Least squares via the normal equations]
This connects to Module~6 of Journey~1 and the preview in Journey~2.
\begin{enumerate}[label=(\alph*)]
  \item When $A\mathbf{x} = \mathbf{b}$ has no solution (overdetermined system), the \textbf{least squares solution} minimises $\|A\mathbf{x} - \mathbf{b}\|^2$. Show the minimum is achieved when $\mathbf{b} - A\hat{\mathbf{x}} \perp C(A)$, i.e.\ when $A^T(A\hat{\mathbf{x}} - \mathbf{b}) = 0$.
  \item Derive the \textbf{normal equations}: $A^T A\hat{\mathbf{x}} = A^T \mathbf{b}$.
  \item When is $A^T A$ invertible? Express the condition in terms of the columns of $A$.
  \item Use the normal equations to fit a line $y = mx + c$ to the data $(1,2), (2,4), (3,5), (4,4)$ by setting up the overdetermined system and solving.
  \item The \textbf{pseudoinverse} $A^+ = (A^T A)^{-1}A^T$ gives $\hat{\mathbf{x}} = A^+\mathbf{b}$. Show that $A^+A = I$ (left inverse) but $AA^+ \neq I$ in general. What does $AA^+$ project onto?
\end{enumerate}
\end{deepproblem}

\begin{deepproblem}[The matrix exponential]
For a square matrix $A$, define $e^A = \sum_{k=0}^\infty \frac{A^k}{k!}$.
\begin{enumerate}[label=(\alph*)]
  \item Show this series converges for any matrix $A$ by bounding $\|A^k\| \leq \|A\|^k$.
  \item Prove $e^{tA}$ solves the matrix ODE $\frac{d}{dt}X = AX$, $X(0) = I$.
  \item If $A = \begin{pmatrix}0 & -1 \\ 1 & 0\end{pmatrix}$, compute $A^2, A^3, A^4$. Use the pattern to sum $e^{tA}$ in closed form.
  \item The system of ODEs $\dot{\mathbf{x}} = A\mathbf{x}$ has solution $\mathbf{x}(t) = e^{tA}\mathbf{x}(0)$. Verify for your matrix $A$ from (c) and initial condition $\mathbf{x}(0) = (1,0)^T$. What geometric motion does the solution trace?
  \item \textit{(Preview.)} If $A$ is diagonalisable with $A = PDP^{-1}$, then $e^A = Pe^DP^{-1}$ where $e^D$ is trivially computed. This makes eigenvalue decomposition (Module~3) essential for solving systems of ODEs.
\end{enumerate}
\end{deepproblem}

\section{Part 6 \texorpdfstring{---}{--} Reflection}

\begin{reflectionbox}
\textbf{Your reflection task.}

\begin{enumerate}
  \item \textbf{On computation.} Row reduction is an algorithm: a fixed sequence of steps that always terminates with the answer. Much of mathematics involves finding algorithms. But not every mathematical question has an algorithm --- some are undecidable. Is linear algebra unusual in being so thoroughly algorithmic? Can you think of questions about matrices that do \emph{not} have efficient algorithms?

  \item \textbf{On the determinant.} The determinant is defined by three axioms and encodes a geometric fact (volume scaling). It is also computed by a specific formula. Write down the three axioms, the geometric meaning, and the $2\times 2$ formula. Do you find it surprising that such a simple formula captures such rich information?

  \item \textbf{On the four subspaces.} The fundamental theorem says the null space and row space are orthogonal complements. This means every vector in $\mathbb{R}^n$ splits uniquely into a component in the row space and a component in the null space. Write a physical analogy: what does this decomposition look like for a specific linear map?

  \item \textbf{Looking ahead.} Module~3 asks: when can we find a basis in which a linear map is diagonal? The matrix $\begin{pmatrix}3&1\\0&2\end{pmatrix}$ from Problem~\ref{prob:change_of_basis} is upper triangular. Is there a basis in which it is diagonal? How would you find such a basis?
\end{enumerate}
\end{reflectionbox}

\medskip
\noindent\textbf{Looking ahead.}
Module~3 is the jewel of linear algebra: eigenvalues and eigenvectors. An eigenvector is a direction that a linear map does not rotate --- only stretches or compresses. Finding eigenvectors means finding the ``natural axes'' of a transformation. This will unify differential equations (the characteristic equation for ODEs is exactly the eigenvalue equation), graph theory (PageRank is an eigenvector), and quantum mechanics (measurable quantities are eigenvalues of operators).

\vspace{1.4cm}
\begin{center}
{\color{MidBlue}\rule{0.5\linewidth}{0.8pt}}\\[6pt]
{\small\color{DarkGrey}\textit{End of Journey 3, Module 2 --- Matrices and Systems}}\\[4pt]
{\small\color{DarkGrey}Continue to: \textbf{Module 3 --- Eigenvalues and Eigenvectors}}\\[6pt]
{\color{MidBlue}\rule{0.5\linewidth}{0.8pt}}
\end{center}

\end{document}